% OWNER: theory-a
% STATUS: review
% SOURCES: notes/SPIS_TRESCI_SZCZEGOLOWY.md §3; GLOSSARY.md; notes/HIPOTEZY_REWIZJA.md
\chapter{Preferencje estetyczne: jawne, ukryte i behawioralne}
\label{ch03:top}
\ChapterStatus{draft}{theory-a}

Rozdział operacjonalizuje preferencje estetyczne jako tworzywo projektu personalizacji w IDA.
Przyjmuje się rozróżnienie kanoniczne: preferencje \textbf{jawne} (\emph{explicit}), \textbf{ukryte} (\emph{implicit}) oraz wskaźniki \textbf{behawioralne}; osobowość (Big Five / IPIP-NEO-120) jest tu jedynie mostem do rozdz.~\ref{ch05:top}.
Modelem docelowym jest synteza wieloźródłowa: sześć kluczy \texttt{GenerationSource} w ślepej macierzy generacji \parencite{idaPlatform2026}.

\section{Preferencje jawne (explicit)}
\label{ch03:sec:explicit}

Preferencje jawne to deklaracje, które użytkownik formułuje świadomie: styl, kolor, materiał, jasność, złożoność.
Klasycznym narzędziem pomiaru znaczenia konotacyjnego jest dyferencjał semantyczny -- pary przeciwstawnych biegunów ocenianych na skali ciągłej \parencite{osgood1957measurement}.
W IDA dyferencjał przyjmuje formę sliderów (m.in.\ ciepło, jasność, złożoność) oraz wyborów palety i materiałów; wynik trafia do pól \texttt{explicit\_*} i dalej do źródła \texttt{explicit} w macierzy.

Jawność nie gwarantuje spójności z zachowaniem wyboru.
Deklaracja „lubię minimalizm” może współistnieć z szybkim akceptowaniem obrazów gęstych wizualnie w teście swipe -- stąd konieczność równoległego kanału ukrytego (rozdz.~\ref{ch03:sec:implicit}) i pytania o rozbieżność (\RQ{2}).

Upload inspiracji (rozdz.~\ref{ch03:sec:inspiracje}) uzupełnia deklarację słowno-skalową o referencję wizualną, nadal jednak w rejestrze jawnym: użytkownik wybiera, \emph{co pokazać} systemowi.

\section{Preferencje ukryte (implicit) i adaptacja IAT}
\label{ch03:sec:implicit}

Test Implicit Association Test (IAT) mierzy siłę skojarzeń poprzez latencje odpowiedzi w zadaniach kategoryzacji, a nie poprzez samoopis \parencite{greenwald1998iat}.
W HCI gest szybkiego wyboru (np.\ swipe) bywa adaptowany jako przybliżenie preferencji trudnych do zwerbalizowania; nie jest to jednak klasyczny IAT laboratoryjny.
\NeedsCite{przeglądy adaptacji IAT / implicit measures w HCI i rekomendacjach -- tylko peer-review z DOI}

W IDA kanał ukryty przyjmuje formę \emph{IAT-like / Tinder}: sekwencja kart z wizualizacjami wnętrz, gest akceptacji lub odrzucenia, zapis w \texttt{participant\_swipes} (m.in.\ \texttt{reaction\_time\_ms}, cechy bodźca).
Agregaty \texttt{implicit\_*} zasilają wariant preferencji ukrytych (\texttt{implicit}) w macierzy.
Ważne zastrzeżenie terminologiczne: „ukryte” oznacza tu pomiar pośredni preferencji estetycznych, a nie tezę psychoanalityczną o nieświadomości.

\begin{figure}[htbp]
  \centering
  \Todo{wstaw plik: figures/ch03/ch03-tinder-iat-flow.pdf --- schemat flow: bodziec wizualny $\rightarrow$ swipe $\rightarrow$ reaction\_time\_ms $\rightarrow$ agregaty implicit\_* $\rightarrow$ źródło implicit}
  \caption{Schemat kanału IAT-like (Tinder) w IDA.
    Opracowanie własne; rama IAT: \textcite{greenwald1998iat}; artefakt: \parencite{idaPlatform2026}.}
  \label{fig:ch03:tinder-iat}
\end{figure}

\section{Rozbieżności implicit versus explicit}
\label{ch03:sec:rozbieznosci}

Model procesów dualnych -- myślenie szybkie i wolne -- dostarcza ramy, w której deklaracja i impulsywna reakcja mogą się rozchodzić \parencite{kahneman2011thinking}.
W personalizacji wnętrz rozbieżność nie jest wyłącznie „szumem pomiarowym”: może wskazywać na konflikt gustu, social desirability albo na trudność werbalizacji atmosfery.

IDA persystuje porównanie kanałów w \texttt{preference\_comparison\_json} (m.in.\ \texttt{style\_match}, zgodność tokenów koloru, biophilia).
Empirycznie pytanie~\RQ{2} bada, czy wyższy mismatch wiąże się z zachowaniem w ślepej macierzy (np.\ dłuższy \texttt{selection\_time\_ms}, częstszy wybór źródeł mieszanych zamiast czystego \texttt{implicit} lub \texttt{explicit}).
Rozbieżność staje się więc dramaturgią poznania gustu -- materiałem projektowym explainable pipeline, a nie błędem do ukrycia.

\section{Wskaźniki behawioralne}
\label{ch03:sec:behavioral}

Obok kierunku swipe platforma zapisuje czas reakcji (\texttt{reaction\_time\_ms}) jako wskaźnik behawioralny przybliżający trudność lub pewność decyzji.
W obecnym pipeline \textbf{nie} są persystowane osobne miary \emph{dwell time} ani \emph{hesitation} jako kolumny badawcze; uproszczenia UI nie należy interpretować jako pełnego śledzenia uwagi.

W macierzy generacji dostępny jest natomiast \texttt{selection\_time\_ms} przy ślepym wyborze -- bliższy decyzji porównawczej niż latencja pojedynczego swipe.
Ograniczenie to zamyka drogę do dawnej hipotezy o wyższości dwell/hesitation nad samym gestem; pozostaje ostrożna analiza czasów reakcji i wyboru (\RQ{2},~\RQ{3}).

\section{Means-end laddering}
\label{ch03:sec:laddering}

Laddering means-end prowadzi od atrybutów produktu (lub przestrzeni), przez konsekwencje użytkowe, do wartości końcowych \parencite{reynolds1988laddering}.
W Room Setup IDA rozmowa (wspierana awatarem przewodniczki) zbiera aktywności, pain points oraz odpowiedzi ladder -- hierarchię „dlaczego to jest ważne”.

Dla generacji obrazu laddering jest mostem między estetyką a funkcją: nie wystarczy wiedzieć, że użytkownik chce „cieplejsze drewno”; trzeba wiedzieć, czy drewno ma wspierać wyciszenie po pracy, koncentrację dziecka, czy rytuał gościnności.
Stąd ścisłe powiązanie z źródłem \texttt{mixed\_functional} oraz z luką PRS current$\rightarrow$target (rozdz.~\ref{ch02:sec:prs}).

\section{Inspiracje wizualne}
\label{ch03:sec:inspiracje}

Upload obrazów referencyjnych wprowadza preferencję przez okaz: użytkownik pokazuje atmosferę, której nie umie w pełni nazwać suwakami.
W produkcji IDA analizę inspiracji i zdjęcia pokoju realizuje warstwa VLM (Gemini Flash Lite); wynik zasila źródło \texttt{inspiration\_reference} w macierzy sześciu wariantów \parencite{idaPlatform2026}.

Eksperymenty w ComfyUI (w tym kontrola referencji i geometrii) stanowią warstwę praktyki autora; produkcyjny artefakt IDA realizuje generację przez Gemini Image / Vertex AI, zachowując inspirację jako jedno z jawnych, cytowalnych wejść briefu -- nie jako niewidzialny styl transferu.

\section{Osobowość: most do rozdziału~\ref{ch05:top}}
\label{ch03:sec:bigfive-bridge}

Model Wielkiej Piątki i instrument IPIP-NEO-120 (domeny + facety) stanowią osobny filar personalizacji -- źródło \texttt{personality} oraz wkład do wariantów mieszanych.
Pełne omówienie konstruktu, mapowań facet$\rightarrow$estetyka, ryzyk determinizmu i etyki wnioskowania o stylu z cech osobowości zawiera rozdział~\ref{ch05:top} (OWNER: theory-b).
Tu odnotowuje się jedynie, że osobowość w IDA jest \emph{hipotezą projektową} podlegającą~\RQ{4}, a nie wyrokiem o „właściwym” guście.

\section{Synteza: model wieloźródłowy i mapa instrumentów}
\label{ch03:sec:synteza}

Preferencje w IDA nie konkurują o miano „prawdziwych”; są kanałami, które explainable pipeline waży i materializuje jako ślepą macierz sześciu źródeł (tab.~\ref{tab:ch03:sources} i~\ref{tab:ch03:instruments}).
Implementacja syntezy i UI macierzy -- rozdz.~\ref{ch07:top}; plan analiz -- rozdz.~\ref{ch06:top}.

\begin{table}[htbp]
  \centering
  \caption{Mapowanie kanałów preferencji na sześć kluczy \texttt{GenerationSource} w IDA.}
  \label{tab:ch03:sources}
  \begin{tabular}{@{}llp{0.42\textwidth}@{}}
    \toprule
    Klucz & Określenie PL & Główne wejścia danych \\
    \midrule
    \texttt{implicit}
      & wariant preferencji ukrytych
      & swipe IAT-like, \texttt{implicit\_*}, \texttt{reaction\_time\_ms} \\
    \texttt{explicit}
      & wariant preferencji jawnych
      & dyferencjał semantyczny, paleta, materiały (\texttt{explicit\_*}) \\
    \texttt{personality}
      & wariant osobowościowy
      & IPIP-NEO-120 / Big Five (rozdz.~\ref{ch05:top}) \\
    \texttt{mixed}
      & wariant mieszany
      & fuzja implicit + explicit (+ wybrane wagi profilu) \\
    \texttt{mixed\_functional}
      & wariant mieszany z funkcją przestrzeni
      & PRS current/target, aktywności, pain points, ladder \\
    \texttt{inspiration\_reference}
      & wariant referencji inspiracji
      & upload inspiracji + analiza VLM \\
    \bottomrule
  \end{tabular}
\end{table}

\begin{table}[htbp]
  \centering
  \caption{Instrumenty preferencji w IDA: źródło naukowe, forma, limity pomiaru.}
  \label{tab:ch03:instruments}
  \begin{tabular}{@{}p{0.18\textwidth}p{0.20\textwidth}p{0.24\textwidth}p{0.26\textwidth}@{}}
    \toprule
    Instrument & Źródło (seed) & Forma w IDA & Limity / uwagi \\
    \midrule
    PRS / mood grid
      & \textcite{hartig1997prs}; ART \parencite{kaplan1995restorative}
      & siatka 2D; \texttt{prs\_ideal/current/target}
      & brak pomiaru post-AI; nie pełne PRS-11 \\
    Biophilia (dawka)
      & \textcite{kellert2008biophilic}
      & 4 poziomy wizualne
      & orientacja wizualna, nie pełna skala biofilii \\
    Dyferencjał semantyczny
      & \textcite{osgood1957measurement}
      & slidery / wybory stylu
      & podatność na social desirability \\
    IAT-like (Tinder)
      & \textcite{greenwald1998iat}
      & swipe; \texttt{participant\_swipes}
      & adaptacja HCI $\neq$ klasyczny IAT; brak dwell/hesitation w DB \\
    Laddering means-end
      & \textcite{reynolds1988laddering}
      & Room Setup + dialog
      & jakość zależna od zaangażowania \\
    Inspiracje
      & praktyka projektowa; VLM w IDA \parencite{idaPlatform2026}
      & upload + analiza
      & bias wyboru referencji przez użytkownika \\
    Big Five / IPIP-NEO-120
      & zob.\ rozdz.~\ref{ch05:top}
      & \texttt{/flow/big-five}
      & most teoretyczny; nie rozwijane tu \\
    \bottomrule
  \end{tabular}
\end{table}

\begin{figure}[htbp]
  \centering
  \Todo{wstaw plik: figures/ch03/ch03-multisource-model.pdf --- diagram: kanały preferencji $\rightarrow$ scoring/brief $\rightarrow$ 6 źródeł $\rightarrow$ ślepy wybór}
  \caption{Model wieloźródłowy preferencji w IDA a macierz generacji.
    Opracowanie własne na podstawie \parencite{idaPlatform2026}.}
  \label{fig:ch03:multisource}
\end{figure}

Tablica źródeł domyka argument rozdziałów~\ref{ch02:top}--\ref{ch03:top}: regeneracyjność i funkcja (PRS, biophilia, ladder) oraz preferencje jawne/ukryte współtworzą materiał, z którego RtD-owy artefakt buduje porównywalne koncepcje wnętrz -- podlegające empirycznej próbie ślepego wyboru (\RQ{3}), a nie autorytetowi pojedynczego kanału gustu.
